% !TeX root = ../report_jouleecosystem.tex
\chapter{Future Work}

Most experiments in this project focused on single-sample inference ($\text{batch size} = 1$), reflecting real-time edge robotics where decisions happen frame-by-frame, before expanding to batch size 32 on the NVIDIA Jetson AGX Orin to study GPU launch amortization. With the end-to-end pipeline now fully established---from physical acquisition to differentiable search and real-silicon deployment (Chapter~\ref{chap:bridge})---several exciting research directions remain open.

\section{Measurement layer (JouleQuest)}

The central open question in the measurement layer is the composability gap: the energy of a full network is close to, but not exactly, the sum of its individual constituent layers. The difference ranges from under $5\%$ on Jetson platforms to roughly $20{-}25\%$ on the Raspberry~Pi~5.

We suspect this divergence stems from cache hierarchy behavior: standalone layer operations execute in tight benchmark loops with hot caches, whereas a complete network rapidly cycles intermediate feature maps through DRAM. Section~\ref{sec:sum-of-parts} outlines the baseline numbers and reasoning. If you pick up this project, the most impactful measurement tasks are:

\begin{itemize}
    \item \textbf{Confirm cache behavior with hardware performance counters (PMUs)}:
    Collect low-level cache-miss metrics using Linux \texttt{perf} or ARM/NVIDIA PMU registers during both standalone and whole-network runs. Correlating last-level cache (LLC) misses and DRAM bus traffic directly with the measured energy gap will turn our cache hypothesis into verified fact.
    
    \item \textbf{Fresh-board validation and inter-board variability}:
    Our original Jetson Nano experienced a hardware failure before final verification. Re-running a core subset of tests on a replacement board will confirm that physical energy profiles remain stable across distinct physical units of the same model.

    \item \textbf{Adaptive and denser lookup grids}:
    JouleGrad estimates unmeasured layer shapes via multilinear interpolation. While accurate along smooth regions, interpolation error grows where hardware execution characteristics change sharply (such as transitions between cache sizes or GPU warp tiling boundaries). Measuring targeted points in these sensitive regions will refine grid precision without requiring an exhaustive brute-force sweep.
\end{itemize}

\section{Estimation and search layers (JouleGrad and JouleNAS)}

With physical validation on real hardware and systematic $\lambda_E$ sweeps now verified, the search framework is ready for broader exploration:

\begin{itemize}
    \item \textbf{Hardware-calibrated composability correction}:
    JouleGrad currently predicts network energy by summing individual layer estimates. Feeding empirical composition data back into the estimator---either through a simple per-board scaling factor $C_{\text{board}}$ or a memory-aware residual term---would allow JouleNAS to optimize directly for composed network energy rather than isolated layer sums.

    \item \textbf{Cross-dataset and transfer learning}:
    Does an energy-optimal mask discovered on CIFAR-10 generalize effectively to other datasets on the same board? Fine-tuning just the linear classifier on a transfer task while freezing the pruned convolutional backbone will reveal whether structural channel sparsity depends on dataset features or primarily reflects underlying hardware efficiency.

    \item \textbf{Extending beyond ResNet backbones}:
    Our search targeted ResNet-18 and SimpleConv. Extending the spec-driven pipeline to depthwise separable networks (MobileNetV3) and edge Vision Transformers (ViTs)---profiling multi-head self-attention and patch embeddings in JouleQuest---will open hardware-aware NAS to the latest generation of edge vision backbones.

    \item \textbf{Dynamic and on-device adaptation}:
    Rather than discovering a single static pruned architecture, the framework could incorporate runtime gating or early-exit branches. Dynamically modulating active channels based on real-time thermal throttling or battery level would bring truly adaptive, hardware-aware execution to edge devices.
\end{itemize}